\documentclass[10pt]{amsart}
\baselineskip=30pt         % double-space
\normalbaselineskip=20pt

\usepackage{amsmath,amssymb,amscd,amsfonts}

\newtheorem{thm}{Theorem}[section]
\newtheorem{lem}[thm]{Lemma}
\newtheorem{rem}[thm]{Remark}
\newtheorem{prop}[thm]{Proposition}
\newtheorem{cor}[thm]{Corollary}
\newtheorem{defn}[thm]{Definition}
\newtheorem{assu}[thm]{Assumption}
\newtheorem{claim}[thm]{Claim}
\newtheorem{ex}[thm]{Example}
\newtheorem{pb}[thm]{Problem}

\newcommand\Pic{{\text{\rm Pic}}}
\newcommand\alb{{ a}}
\newcommand\lra{\longrightarrow}
\newcommand\Alb{{ A}}
\newcommand\Mod{\text{\rm Mod}}
\newcommand\ot{{\otimes}}
\newcommand\OO{{\mathcal{O}}}
\newcommand\HH{{\mathcal{H}}}
\newcommand\QQQ{{\mathcal{Q}}}
\newcommand\PP{{\mathcal{L}}}
\newcommand\LL{{\mathcal{P}}}
\newcommand\FF{{\mathcal{F}}}
\newcommand\GG{{\mathcal{G}}}
\newcommand\PPP{{\mathbb{P}}}
\newcommand\Z{{\mathbb{Z}}}
\newcommand\Q{{\mathbb{Q}}}
\newcommand\QQ{{\mathbb{Q}}}
\newcommand\CC{{\mathbb{C}}}
 \newcommand\C{{\mathbb{C}}}
 \newcommand\II{{\mathcal{I}}}
\newcommand\Hom{{\mathcal{H}om}}
 \newcommand\sh{{\hat {\mathcal S}}}
 \newcommand\s{{ {\mathcal S}}}
%%%%%%%%%%%%%%%%%%%%%%%%%
%
%   Lecture plan
% L1=1,2 
% L2=3,4
% L3=5,6
% L4=7
% L5=8
%
%
%
%
%%%%%%%%%%%%%%%%%%%%%%%%%
 \title{On the birational geometry of Irregular Varieties}
 \author {Christopher D. Hacon}
 \date{}
 \begin{document}
 \abstract{The purpose of this course is to illustrate how the theory of Fourier
Mukai transforms can be used to give surprisingly precise results
in the birational classification of irregular varieties.}
 \endabstract
 \maketitle
 %%%%%%%%%%%%%%%%%%%%%%%%%%%%%%%%%%
 \section{Introduction}
 %%%%%%%%%%%%%%%%%%%%%%%%%%%%%%%%%%
Let $X$ be a smooth complex projective variety
and $\omega _X$ the canonical
 line bundle. We will denote by
\begin{itemize}
\item $P_m(X):= h^0(\omega _X^{\ot m})$ the $m$-th {\em plurigenus of} $X$.
\item $q(X):=h^0(\Omega ^1_X)$ the {\em irregularity} of $X$.
\end{itemize}
These numbers are birational invariants of the variety $X$. (Recall that
two varieties $X,X'$ are birational if and only if they have 
some isomorphic open subset $U\cong U'$.)
They are interesting because when they are positive, $X$ comes
equipped with a natural algebraic map.

If $P_m(X)>0$, then the sections of $H^0(X,\omega _X^{\ot m})$  
define a rational map
$$\varphi ^m_X:X \dashrightarrow \PPP \left( 
H^0(X,\omega _X ^{\ot m})\right) .$$
Recall that the {\em Kodaira dimension}  
$$\kappa (X)\in \{-\infty,0,1,...,\dim (X)\},$$ 
is defined as follows: if $P_m(X)=0$ for all
$m>0$ we set $\kappa (X)=-\infty$, and otherwise $\kappa (X)$ is 
the maximum of the dimensions of the images
of $\varphi ^m_X$ for all $m>0$. This is also a birational invariant of $X$.

If $q(X)>0$, then by integrating the holomorphic one forms
$\omega \in H^0(X,\Omega ^1_X)$, one obtains a morphism
$$\alb :X \lra \Alb (X):= \frac {H^0(X,\Omega ^1_X)^\vee }{H_1(X,\Z )}.$$
The dimension of the image of $\alb$ is the {\em Albanese dimension of} $X$.
The Albanese variety $\Alb (X)$ is an abelian variety i.e. a complex 
projective torus. Its dual abelian variety 
$$\widehat {\Alb (X)}=\Pic ^0(\Alb )\cong\Pic ^0(X)$$
parametrises the topologically trivial line bundles on $X$ and $\Alb$.

One of the standard questions in Algebraic geometry is to study the 
geometry of these natural maps and their interaction.
For example one may ask:

\bigskip \noindent {\bf Question:}
{\it If $X$ is of general type (i.e. $\kappa (X)=\dim (X))$, 
then what is the smallest
value of $m$ for which $\varphi^m_X$ is a birational map?}
\bigskip

If $\dim (X)=1,2$ then $m=3,5$ will suffice. However, when $\dim (X)\geq 3$
any explicit estimate seems extremely hard to obtain. 
In contrast to this, it is shown in \cite{CH3} 
that for varieties with maximal Albanese dimension, $m=6$ always 
suffices (regardless of the dimension of $X$!).  




Ideally, one could even hope to characterize (the birational equivalence class
of) a projective variety $X$ by its birational invariants. 
For example recall the following result for surfaces:

\bigskip \noindent 
{\bf Castelnuovo's Criterion:} {\it If $P_2(X)=q(X)=0$ then $X$ is a 
rational surface.}

\bigskip \noindent
{\bf Enriques:} {\it If $P_2(X)=1$ and $q(X)=2$, then $X$ is birational to an 
abelian variety.}

\bigskip
Classification results of this kind are extremely rare, 
however there is a famous 
result of Kawamata that generalizes Enriques' result to higher dimensions:
\begin{thm} \label{KAW} \cite{Ka1}
If $X$ is a smooth complex 
projective variety with $\kappa (X)=0$ then 
its Albanese morphism is surjective (in particular $q(X)\leq \dim (X)$).
If also $q(X)=\dim (X)$, then $X$ is 
birational to an abelian variety.
\end{thm} 
This 
is somewhat unsatisfactory as the Kodaira dimension is
not a naive invariant, it depends on $P_m(X)$ for all $m>0$.
Koll\'ar has proved an effective version of this result: {\em If $X$ is
a smooth complex projective variety with $P_m(X)=1$ for some $m\geq 3$,
then its Albanese map is surjective. If $P_m(X)=1$ for some $m\geq 4$
and $q(X)=\dim (X)$, then $X$ is birational to an abelian variety.}
In \cite{CH1}, Koll\'ar's result was refined as follows:

\bigskip\noindent
{\bf Chen-Hacon}: 
{\it If $P_m(X)=1$, for some $m\geq 2$, then the Albanese map
$\alb :X \lra \Alb$ is surjective. If also $\dim (X)=q(X)$ then 
$X$ is birational to an abelian variety.}
\bigskip

Recently it has been shown that if $P_3(X)=2$ (resp. $P_3(X)=3$) 
and $q(X)=\dim (X)$, then $X$ is a double (resp. bi-double) cover of
an abelian variety cf. \cite{HP1}, \cite{Hac2}.

\bigskip
The main techniques used in obtaining these kinds of results are:

\begin{itemize}

\item The theory of Fourier-Mukai transforms \cite{Mu} which characterizes 
coherent sheaves $\FF$ on an abelian variety $A$, 
in terms of the cohomology groups $H^i(A,\FF \ot P)$ with
$P\in \Pic ^0(A)$. (Actually a sheaf theoretic version of this 
statement is necessary.)

\item The Generic Vanishing Theorems of Green and Lazarsfeld, 
which for example imply
that if $X$ is a smooth projective variety 
of maximal Albanese dimension, then $\chi (\omega _X)\geq 0$.

\item The vanishing theorems of Nadel, Kawamata, Viehweg and the 
results on higher 
direct images of dualizing sheaves due to Koll\'ar.
\end{itemize}
 
\newpage


 %%%%%%%%%%%%%%%%%%%%%%%%%%%%%%%%%%
 \section{Sheaves on abelian varieties}
 %%%%%%%%%%%%%%%%%%%%%%%%%%%%%%%%%%
 Let $A$ be a $g$-dimensional abelian variety and $\hat{A}=\Pic ^0(A)$ its 
dual abelian variety. We will denote by $p_A,p_{\hat {A}}$ the projections
of $A\times \hat{A}$ on to $A,\hat {A}$. Let $\LL $ be the 
normalized Poincar\'e 
line bundle on $A\times \hat {A}$.
So for $\hat {a}\in \hat {A}$, the line bundle $\LL |A\times \hat {a}$
is just the topologically trivial line bundle $\LL _{\hat{a}}$
corresponding to the
point $\hat {a}\in \Pic ^0(A)=\hat{A}$.
The main tool in understanding sheaves on an abelian variety is the 
Theory of Fourier-Mukai transforms. Recall that Mukai defines the 
functor $\s _A :D^b_{qc}(A)\lra D^b_{qc}(\hat {A})$ by
$$\s_A (F)=Rp _{\hat{A},*}(\LL \ot ^L p_A ^*M).$$ 
Similarly 
one defines $\s _{\hat{A}}:D^b_{qc}(\hat {A})\lra D^b_{qc}(A)$.
Here $D^b_{qc}(A)$ denotes the derived category of quasi-coherent sheaves
on $A$ with bounded cohomology. We will mostly just be interested
in the case where $F$ is a sheaf. In this case, for general
$\hat {a} \in \hat {A}$, one has that $H^i(\s _A (F))\ot \CC (\hat {a})
\cong H^i( F\ot \LL _{\hat{a}})$.
 By \cite{Mu}, one has: 
 \begin{thm} \label{mukai} {\bf (Mukai)} There exists isomorphisms of functors
 $$\s _A \circ \s _{\hat {A}}\cong (-1_{\hat{A}})^*[-g]$$
 and
 $$\s _{\hat {A}}\circ \s _A \cong (-1_{{A}})^*[-g].$$
 \end{thm}
Here $[-g]$ means: shift the corresponding complex $g$ spaces to the right.
Loosely speaking this means that one can recover a sheaf $F$ from its 
cohomology groups $H^i(F\ot P)$. 

\begin{ex}\label{E1} If 
$H^i(A,F\ot P)=0$ for $0\leq i \leq g$ and all $P\in \Pic ^0(A)$,
then by $\s _A(F)=0$ and so $F=0$. 
Similarly, if $\phi :F\lra G$ is an inclusion of sheaves
inducing isomorphisms $H^i(F\ot P)\lra H^i(G\ot P)$ for $0\leq i \leq g$ 
and all $P\in \Pic ^0(A)$, then $F\cong G$.
\end{ex}
\begin{ex}\label{E2} Let $F=\CC (\hat {a})$ for some point $\hat{a}\in 
\hat{A}$, then $S_{\hat{A}}(F)=\LL _{\hat {a}}$ (i.e. the complex
given by $F^i=0$ for $i\ne 0$ and $F^0=\LL _{\hat {a}}$).
Therefore, $S_A(\LL _{\hat {a}})=\CC (\hat {a})[-g]$, in particular
$R^ip_{\hat{A},*}(p_A^* \OO _A \ot \LL)=0$ for $i\ne g$ and
$R^gp_{\hat{A},*}(p_A^* \OO _A \ot \LL)=\CC (0_{\hat{A}})$. 
\end{ex}


 \begin{ex} \label {exM} {\bf (Mukai)}
 A vector bundle $U$ on $A$ is unipotent if there exists a filtration
 $$0=U_0\subset U_1\subset ...\subset U_{n-1}\subset U_n=U$$
 such that $U_i/U_{i-1}\cong \OO _A$ for all $1\leq i\leq n$. One has that
 $R\sh (U)=R^g\sh (U)$ and $R^g\sh (U)$ is supported on
 $\hat {0}\in \hat{A}$. This gives an equivalence between the category of
 unipotent vector bundles on $A$ and the category of coherent sheaves on
 $\hat {A}$ supported on $\hat{0}$ i.e. the category of Artinian $\OO 
 _{\hat{A},\hat{0}}$-modules.
 \end{ex}
\begin{ex} \label{Ex4}
Let $\FF$ be a coherent sheaf on $A$ such that for all $P\in
\Pic ^0(A)$ one has
$$h^0(\FF \ot P)=1,\ \ and\ \ h^i(\FF \ot P)=0\ \ \forall\ i>0.$$
Then by cohomology and base change $L:= \s _A (\FF )=H^0
\s _A (\FF )$ is a line bundle on $\hat {A}$.
Moreover, one has that $\s _{\hat{A}}(L)=H^g\s _{\hat {A}} (L)$ and
hence $h^g(L\ot P_x)\ne 0$ for some $x\in A$. 
Line bundles on an abelian variety are particularly well understood.
One sees that $L$
is negative semidefinite (i.e. $L^\vee$ is nef). 
It follows that there exists a (surjective with connected fibers)
homomorphism of abelian varieties $p:\hat{A}\lra B$
and a non-degenerate line bundle $M$ on $B$ such that $L\ot P_x=p^*M$.
In particular $h^g(L\ot P_x \ot Q)=0$ for all $Q\notin \hat {B}$
and therefore, $\FF$ is supported on the abelian subvariety $\hat{B}
\hookrightarrow A$.

If the generic rank of $\FF$ is 1, then $h^g(L \ot P)=1$ for general
$P\in \Pic ^0(\hat{A})=A$ and so $L$ is dual to a principal polarization
and hence $\FF$ is a principal polarization (and in particular
a line bundle).

\end{ex}

% We also recall that 
% $$D_{\hat {A}}\circ \s _A \cong ((-1_{\hat {A}})^*\circ 
% \s _A \circ D_A)[g].$$
%(For an $n$-dimensional variety $X$, 
%the dualizing functor is defined by $D_X(?)=
%R\HH om (?, \omega ^._X[n])$ where $\omega ^._X$ is the dualizing
%complex of $X$. Therefore, $D_A= R\HH om (?, \OO _A[g])$.)
\newpage
%%%%%%%%%%%%%%%%%%%%%%
 \section{Generic Vanishing Theorems}
%%%%%%%%%%%%%%%%%%%%%%%%%%%%%%%%%%
For $X$ a smooth complex projective variety (or more generally a K\"ahler
manifold), one can define {\em cohomological support loci}  
$$V^i(\omega _X):=\{ P\in \Pic ^0(X)| h^i(\omega _X \ot P)\ne 0\}.$$
In view of the Fourier-Mukai transform, one would expect the $V^i(\omega _X)$
to ``determine'' the sheaves $\alb _* (\omega _X)$, so
it is not surprising that
these loci are very important in the study of irregular varieties.
Their geometry is governed by the following result of Green and Lazarsfeld
\cite{GL}
\begin{thm}\label{GVT}(Generic Vanishing Theorem) 

(1) Any irreducible component of $V^i(X,\omega _X )$ is a translate\footnote 
{If $X$ is 
(deformation equivalent to)
a projective manifold, then Simpson has shown that: {\it 
Every irreducible component of $V^i(X,\omega _X)$ is
a torsion translate of an abelian  subvariety of $\Pic ^0(X)$.}} of
an abelian subvariety of $\Pic ^0(X)$
and is of codimension at least
$i-(\dim (X) -\dim (\alb _X(X)))$.

(2) Let $P \in T$ be a general point of an irreducible component
$T$ of $V^i(X,\omega _X)$.
Suppose that $v
\in H^1(X, \OO
_X)\cong T_{P }\Pic ^0(X)$ is not tangent to $T$. Then the sequence
$$H^{i-1}(X,\omega _X \otimes P) \overset {\cup v}  {\longrightarrow}
H^{i}(X,\omega _X \otimes P)  \overset {\cup v}  {\longrightarrow }
H^{i+1}(X, \omega _X \otimes P) $$
is exact.
If $v$ is tangent to $T$, then the maps in the above sequence vanish.

(3) If $X$ is a variety of maximal Albanese dimension, then
$$\Pic ^0(X)\supset V^0(X,\omega _X )\supset V^1(X,\omega _X)
\supset \ ...\ \supset V^n(X,\omega _X )=\{ \OO _X\}.$$

\end{thm}
The above result is even true for K\"ahler manifolds and its proof
depends on Hodge theory. 
The idea behind this result is that the cohomology of the 
above complex measures
the first order deformations of cohomology classes in
$H^i(X,\omega _X)$, and that if a class deforms to first order, then
all higher obstructions automatically vanish.

From (1), one sees that if $X$ has maximal Albanese dimension, then
$\chi (\omega _X)\geq 0$. If $X$ is a subvariety of general type
of an abelian variety, then it has been shown that $\chi (\omega _X)>0$.
It was conjectured by Koll\'ar that 
(as in the case of surfaces of general type),
if $X$ is of  
maximal Albanese dimension and general type, then $\chi (\omega _X)>0$
and $P_1(X)\geq 2$. This is however not the case (cf. \cite{EL}, \cite{CH4}).

Green and Lazarsfeld also conjecture the following statement which is
closely related to their generic vanishing theorems (cf. \cite{GL}):

\bigskip \noindent{\bf Conjecture:} {\em Let $X$ be a smooth complex projective
variety with maximal Albanese dimension, then for the universal family
of topologically trivial line bundles $\PP\lra X\times \Pic ^0(X)$,
one has that $$R^i\pi _{\Pic ^0(X),*}(\PP)=0$$
for all $i<\dim (X)$.}

In section \ref{GL}, 
we will show that this is true, and if $X$ is not of 
maximal Albanese dimension, an appropriate statement still follows.


\newpage
%%%%%%%%%%%%%%%%%%%%%%%%%%%%%%%%%%%%%%%
\section{\label{DC}A derived category approach to generic vanishing theorems}
 %%%%%%%%%%%%%%%%%%%%%%%%%%%%%%%%%%%%%%%
Let $L$ be an ample vector bundle on $\hat{A}$. 
By \cite{Mu} Proposition 3.11, 
 $\hat {L}:=H^0\s _A(L)\cong  \s _A(L)$ is a vector bundle on $A$ of rank
 $h^0(L)$. Moreover,
 $$\phi _L^* \left( \hat{L}^\vee \right) 
\cong \left( L \right) ^{\oplus h^0(L)}.$$
(Recall that the isogeny $\phi _L:\hat {A}\lra A$ is defined by $\phi _L
(\hat {a})=t^*_{\hat {a}}L^\vee \ot L$.)

 \begin{thm} \label{TDC} \cite{Hac4}
 Let $\FF$ be a coherent sheaf on an abelian variety $A$.  
 The following are equivalent:
 \begin{enumerate}
 \item 
 For any sufficiently 
 ample line bundle
 $L$ on $\hat{A}$,  
 $$H^i(A, \FF \ot \hat {L}^\vee )=0\ \ \ \forall \ i>0;$$
 \item 
 There is an isomorphism
 $$\s _A D_A(\FF ) \cong H^0 \left( \s _A
 D_A(\FF )\right) .$$
 \end{enumerate}
 \end{thm}
Recall that for an $n$-dimensional variety $X$, 
the dualizing functor is defined by $D_X(?)=
R\HH om (?, \omega ^._X[n])$ where $\omega ^._X$ is the dualizing
complex of $X$. Therefore, $D_A= R\HH om (?, \OO _A[g])$.

 \begin{proof}
 By Grothendieck Duality (G.D.)\footnote{(G.D.) 
Let $f:X\lra Y$ be a morphism of 
projective varieties, $F\in D^b_{qc}(X)$, then $Rf_*D_X(F)=D_Y(Rf_*(F))$.
In particular one has $R\Gamma (D_XF)\cong D_k(R\Gamma (F))$ where $R\Gamma 
=R\gamma _*$ and $\gamma : X \lra Spec (k)$ is the structure morphism.}
and the projection formula 
(P.F.)\footnote{(P.F.) Let $f:X\lra Y$ be a proper morphism of 
quasi-projective varieties, then there is an isomorphism
$Rf_*(F)\ot G\lra Rf_*(F\ot Lf^*G)$.} it follows that
$$ D_k (R \Gamma (\FF \ot 
 \hat {L}^\vee ))\cong^{G.D.}R \Gamma (D_A(\FF \ot \hat{L}^\vee ))\cong$$ 
 $$R \Gamma (D_A(\FF )\ot \hat{L})\cong 
R \Gamma (D_A(\FF )\ot p_{A,*}(\LL \ot p_{\hat{A}}^*L))\cong ^{P.F.}$$
 $$R \Gamma (L p_A^* D_A (\FF) \ot \LL \ot p_{\hat{A}}^*L)\cong^{P.F.}
R \Gamma (\s _A (D_A(\FF ))\ot L).$$
In particular, $D_k (R \Gamma (\FF \ot 
 \hat {L}^\vee ))$ is a sheaf if and only if $R \Gamma (\s _A
 (D_A (\FF)) \ot L)$ is a sheaf.
The theorem now follows from the following remarks:
\begin{enumerate} 
\item Notice that there is an isomorphism
$$H^0(A,\FF \ot \hat {L}^\vee )^\vee \cong D_k (R \Gamma (\FF \ot 
 \hat {L}^\vee ))$$ 
if and only if $$H^i(A,\FF \ot \hat {L}^\vee )=0\ for\ all\ i>0.$$
 
\item For $L$ sufficiently ample, the coherent sheaves
 $H^j\s _A (D_A (\FF))\ot L$ are globally generated and 
have vanishing higher cohomology. 
Consider the spectral sequence
 $$E_2^{i,j}=R^i\Gamma (H^j\s _A (D_A (\FF ))\ot L)\Rightarrow R^{i+j}\Gamma (
 \s _A (D_A(\FF )) \ot L).$$
 Since $E_2^{i,j}=0$ for all $i\ne 0$, this sequence degenerates at the $E_2$
 level and hence $$H^0(\hat{A}, H^j\s _A (D_A(\FF ))\ot L)\cong 
 R^j\Gamma (\s _A ( D_A(\FF ))\ot L ).$$
Therefore, 
$$R \Gamma (\s _A ( D_A (\FF))\ot L)\cong
R^0 \Gamma ( \s _A (D_A (\FF)) \ot L)$$
if and only if 
$$ \s _A (D_A(\FF ))\cong H^0\s _A (D_A(\FF )).$$
\end{enumerate}
 \end{proof}
\begin{cor} \label{C} Let $\FF$ be a coherent sheaf as above. For all $i>0$, 
the support of $H^i\s _A (\FF )$ 
has codimension at least $i$ in $\hat {A}$.
\end{cor}
\begin{proof} Since $D_{\hat {A}}\circ \s _A \cong (-1_A)^*\circ \s_{\hat{A}}
\circ D_A [g]$ (cf. \cite{Mu} (3.8)), one has that
$$\s _A (\FF) \cong 
D_{\hat{A}} \left((-1_{\hat {A}}^*\circ  \s _A 
\circ D_A)(\FF)[g]\right)\cong
R\Hom \left( \GG,\OO _{\hat{A}} \right),$$
where $\GG:=-1^*_{\hat{A}}( \s _A ( D_A(\FF )))$ is a sheaf.
It suffices therefore to show that for any coherent sheaf $\GG$ on a smooth 
affine variety $Y=Spec(A)$, one has that
the sheaves $R^i\Hom( \GG, \OO _Y)$ are supported in codimension at 
least $i$.
Let $M=\Gamma (Y,\GG )$ and $W$ be an irreducible component of the support of 
$R^i\Hom ( \GG, \OO _Y)={Ext^i(M,A)}^\backsim$. Let $P\in Spec (A)$ such that
$W=Spec (A/P)$. 
Since localization is exact, one has that
$$0\neq Ext^i (M,A)\ot   A_P\cong
Ext^i (M_P,A_P).$$
Since $A$ is regular, it follows that $A_P$ is regular and hence that
$$i\leq\dim A_P=\dim (Y)-\dim (W)$$ (cf. \cite{Ha1} \S III.6).
\end{proof}
 \begin{cor}\label{C1} Let $\FF$ be a coherent sheaf as above, 
and $P\in \hat{A}
$,
 then:
 \begin{enumerate}
 
 \item If $i\geq 0$ and $H^i(A,\FF \ot P)=0$, then $H^{i+1}(A,\FF \ot P)=0$.

\item For all $i>0$ any irreducible component of the loci
$$V^i(\FF ):=\{ P\in \Pic ^0(A)\ |\ h^i(\FF \ot P)\ne 0\}$$
has codimension at least $i$ in $\hat {A}$.
 
 \item If $H^0(A, \FF \ot P)=0$ then $H^0\s _A (D_A(\FF))\ot k(P^\vee )=0$.

 \item If $H^0(A,\FF \ot P)=0$ for all $P\ne \OO _A$, then $\FF$ is a 
 unipotent vector bundle.
 \end{enumerate}
 \end{cor}
 \begin{proof} 
 (1) Since
$$0=H^i(A,\FF \ot P)^\vee \cong H^{-i}(D_\CC R
 \Gamma (\FF \ot P))\cong R^{-i}\Gamma (D_A(\FF )\ot P^\vee ),$$
by Cohomology and Base Change, one sees that the natural
homomorphism
$$0=H^{-i-1}\s _A (D_A(\FF ))\ot k(P^\vee)
\lra R^{-i-1}\Gamma (D_A(\FF)\ot P^\vee)
$$ 
is surjective and hence that $H^{i+1}(A,\FF \ot P)=0$.

(2) Let $W$ be an irreducible component of $V^i(\FF)$ and $P$ 
a general point in
 $W$.
 Let $j$ be the largest integer such that $W\subset V^j(\FF )$.
By Cohomology and Base Change (for sheaves), one 
has that $H^j\s _A (\FF )\ot k(P)
\ne 0$
and by the previous corollary, one has that $W$ has codimension at least $j$.

 (3) This follows immediately from Cohomology and Base Change.

 (4) By (3), we have that $H^0\s _A (D_A(\FF ))$ is supported on 
 the closed point $\hat{0}\in \hat {A}$ corresponding to $\OO _A$.
 By Example \ref{exM}, $H^0 \s _A (D_A(\FF ))$ is an Artinian
 module and hence
 $$D_A(\FF )\cong (-1_A)^* \s _{\hat {A}}  \s _A (D_A (\FF ))[-g]$$
 is (the shift of) a unipotent vector bundle say $V$. Finally it is easy to see
 that $$\FF=D_A(D_A(\FF))=V^\vee$$ is also a unipotent vector bundle.
 \end{proof}

\newpage
%%%%%%%%%%%%%%%%%%%%%%%%%
\section{\label{GL}A conjecture of Green and Lazarsfeld}
%%%%%%%%%%%%%%%%%%%%%%%%% 

We will need the following results about higher direct images of dualizing
sheaves:

 \begin{thm} \label{kollar} {(\cite {Ko1}, \cite{Ko2})}
 Let $X,Y$ be complex projective varieties of dimension $n,k$
 with $X$ smooth.
 Let $f:X\lra Y$ be a surjective map and $L$ an ample line bundle on $Y$. Then
 \begin{enumerate}
 \item $R^if_* \omega _X $ is torsion free for $i\geq 0$,
 \item $H^j(Y,L\ot R^if_*\omega _X )=0$ for $j>0$,
 \item $R f_* \omega _X \cong \sum R^if_* \omega _X [-i],$
 \item $R f_* \OO _X \cong \sum {}^dR^if_* \OO _X$, where 
 ${}^dR^if_* \OO _X=D_Y(R^{n-k-i}f_* \omega _X [k+i])$.
 \end{enumerate}
 \end{thm}
 
For future reference, we also recall the following vanishing theorem 
(cf. \cite{Ko3} \S 10).
\begin{thm} \label{omegavan}Let $f\colon X\to  Y$ be a  surjective 
map of projective
varieties, $X$ smooth, $Y$ normal. Let $M$ be a line bundle on $X$
such that $M\equiv N+f^* L +\Delta$, where $N$ is a 
$\Q$-divisor on $X$ which is either
nef and big or numerically trivial, $L$
is a $\Q-$divisor on $Y$
and $\Delta$ has normal crossing support with 
$\lfloor \Delta \rfloor=0$. Then
\begin{itemize}
\item $R^jf_*(\omega _X\ot M) $ is torsion free for $j\geq 0$;

\item If $L$ is nef and big, and $g:Y\lra Z$ is any morphism with $Z$ 
projective,
then $H^i(Y,R^jf_*(\omega _X\ot M))
=0$ for $i>0$, $j\geq 0$ and
$$H^i(Y,R^jf_*(K_X+M))=H^i(Z,g_*(R^jf_*(K_X+M))),\ \ i,j\geq 0.$$


\end{itemize}
\end{thm}

 When $X$ is a smooth complex projective variety and $\alb :X\lra A$ is the
 Albanese morphism, the results of the previous section
allow us to recover a generalization of the conjecture of
 Green and Lazarsfeld mentioned above.
To fix the notation, let $\PP=(\alb ,id_{\hat{A}})^*\LL$ and 
let $\pi _{\hat A}$ denote
 the projection of $X\times \hat {A}$ onto $\hat{A}$.
 \begin{thm} \label{conjGL} \cite{Hac4}
 Let $X$ be a smooth complex projective variety of dimension 
 $n$, and Albanese dimension $k:=\dim \alb (X)$. 
 Then $$R \pi _{\hat {A },*}(\PP )\cong \sum _{i=0}^{n-k}H^{k+i}
 \s _A ({}^dR^i\alb _* \OO _X)[-k-i].$$
 In particular, when $X$ is of maximal Albanese dimension (i.e. $k=n$),
 one has that $R \pi _{\hat {A},*}(\PP )\cong 
H^n\s _A ({}^d R^0\alb_* \OO _X)$ is
 a sheaf.
 \end{thm}
\begin{proof} 
 Let $Y:=\alb (X)$. We have that
 $$R \pi _{\hat {A},*}(\PP )\cong R \pi _{\hat {A},*}
((\alb \times id_{\hat{A}})
^*
 \LL)\cong ^{P.F.}
 R p_{\hat {A},*}( L p_A^* (R \alb _* \OO _X)\ot \LL )\cong$$
 $$R p_{\hat {A},*}( L p_A^*
 (\sum _{i=0}^{n-k}{}^dR^i\alb _* \OO _X)\ot \LL)
 \cong 
 \sum _{i=0}^{n-k} \s _A ({}^dR^i\alb_* \OO _X).$$
 Since $${}^dR^i\alb _* \OO _X\cong D_Y(R^{n-k-i}\alb _* \omega _X[k+i])\cong
 D_A(R^{n-k-i}\alb _* \omega _X[k+i])$$
 it suffices to show that the sheaves $\FF ^{n-k-i}:= 
R^{n-k-i}\alb _* \omega _X
$
 satisfy the hypothesis of Theorem \ref{TDC}. 
 This is however an immediate consequence of Theorem \ref{kollar} and the  
observations that follow:

 Let $\FF =\FF ^{n-k-i}$. As mentioned in the section \ref{DC},
 $\hat {L}:=H^0\s _A(L)\cong  \s _A(L)$ is a vector bundle on $A$ of rank
 $h^0(L)$ and
 $$\phi _L^* \left( \hat{L}^\vee \right) \cong \left( L \right) 
^{\oplus h^0(L)}
.$$
 We have that 
 $H^i(A,\FF \ot \hat {L}^\vee )$ is a direct summand of 
 $$H^i(A,\phi _{L,*}(\OO _{\hat {A}})\ot \FF \ot \hat {L}^\vee  )
\cong ^{P.F.}H^
i(\hat {A},
 \phi _L^* (\FF \ot \hat {L}^\vee ))\cong \bigoplus _{i=1}
^{h^0(L)}H^i(\hat {A},
 \phi _L^* (\FF )\ot L ).$$ So it suffices to show that $H^i(\hat {A},
 \phi _L^* (\FF )\ot L )=0$ for all $i>0$.
 Let $X^\prime:=X\times _A \hat {A}$ and $\alb ^\prime:X^\prime
\lra \hat {A}$ be the induced morphism.
 By flat base change, $\phi _L^* (\FF )\cong R^{n-k-i}{\alb }^\prime_* \omega 
_{X^\prime}$, and so by
 Theorem \ref{kollar}, we have the required vanishing: $$h^i(\hat {A}, 
 R^{n-k-i}{\alb }^\prime_* \omega _{X^\prime}\ot L)=0\ for \ all\ i>0.$$

 \end{proof}
 \begin{cor} \label{C2}
Let $\alb:X\lra \Alb$ be as in the preceding theorem, then
\begin{enumerate}
\item Every irreducible component of the loci $V^i(R^ja_*\omega _X)$ has
codimension at least $i$ in $\hat {A}$.

 \item  For all $i<\dim \alb (X)$ and general $P\in \Pic ^0(X)$,
 one has that $H^i(X,P)=0$.

 \item If $H^j(R^i\alb _*\omega _X\ot P)=0$, then 
$H^l(R^i\alb _*\omega _X\ot P)
=0$
 for all $l>j\geq 0$ and $0\leq i\leq k$.
\item If $H^0( R^i\alb _*\omega _X\ot P)=0$ for all $P\ne \OO _A$,
then $ R^i\alb _*\omega _X$ is a unipotent vector bundle.
\end{enumerate} 
\end{cor}
 \begin{proof} This is analogous to Corollary \ref{C1}.
 \end{proof}
 \begin{cor} \label{k0}
If $\kappa (X)=0$, then the Albanese map $\alb :X \lra A$
is surjective.
 \end{cor}
 \begin{proof} By a result of Fujita,
there is a generically finite cover $\tilde {X}\lra X$ 
with $\kappa (\tilde {X}
)=0$ and $P_1(\tilde {X})=1$. We claim that since $\kappa (\tilde {X}
)=0$, then $h^0(K_{\tilde {X}}\ot P)=0$ for all $P\ne \OO _{\tilde {X}}$.
If this where not the case, then by Simpson's resulf (cf. Theorem
\ref{GVT}), one can assume that $h^0(K_{\tilde {X}}\ot P)\ne 0$
for some torsion element $P\in \Pic ^0(X)$. Pick a divisor 
$D\in |K_X+P|$, then for sufficiently divisible $t$, one has
that $tD\in |t(K_X+P)|=|tK_X|$ and hence $h^0(2K_X)>1$ which is
impossible. 
By the above corollary, $\alb _{\tilde {X},*}\omega _{\tilde {X}}$ is
a unipotent vector bundle on $\Alb (\tilde {X})$ and hence
$ {\alb _{\tilde {X}}}:\tilde {X}\lra \Alb (\tilde {X})$ is surjective.
It follows easily that $X\lra A(X)$ is also surjective.
  \end{proof}
%\begin{cor} Let $f:X\lra Y$ be a morphism of smooth projective varieties
%which induces an inclusion $\omega _Y\lra f_* \omega _X$.
%If $h^0(\omega _Y\ot P)=h^0(f_* \omega _X\ot P)$ for all $P\in 
%\Pic ^0(Y)$, then $\alb _{Y,*}(\omega _Y)=\alb _{Y,*}(\omega _X)$.
%\end{cor}
% \begin{proof} Let $K=coker (\omega _Y\lra f_* \omega _X)$
% \end{proof}

\newpage 
%%%%%%%%%%%%%%%%%%%%%%%%%%%%%%%%%%%%%%%%%%%
\section{Effective criteria for birational morphisms}
%%%%%%%%%%%%%%%%%%%%%%%
Throughout this section, $f:X\lra Y$ denotes a surjective
morphism of smooth complex projective varieties, with
$\dim (X)=\dim (Y)=n$. Recall that $X$ is of {\em general type} if 
$\kappa (X)=\dim (X)$.
\begin{thm}\label{birational} {\bf (Hacon-Pardini)} \cite{HP1}
 If $Y$ is of  general type and maximal Albanese dimension,
then $f:X\lra Y$ is birational if and only if $P_m(X)=P_m(Y)$ for some
$m\geq 2$.
\end{thm}
\begin{proof} Assume for simplicity that $m=2$. The Albanese morphism
$\alb _{Y}:Y\lra \Alb (Y)$ is generically finite onto its image. 
We would like to compare
$\alb_{Y,*}(f_* \omega _X ^{\ot 2})$ and $\alb _{Y,*}(\omega _Y^{\ot 2})$. If
these are generically isomorphic, then the generic degree
of $f$ is 1. The first step is to extract the positive part of 
$\omega _Y,\omega _X$.

Let $H$ be any ample divisor on $\Alb (Y)$, for some $t\gg 0$,
one has that $|tK_Y-\alb _Y^*H|$ is nonempty.
Fix a general divisor $$B_Y\in |tK_Y -\alb _Y^*H|$$ and let
$B_X:=f^*B_Y+tK_{X/Y}$.  After replacing $X,Y$ by appropriate birational 
models, we may  assume that $B_X,B_Y$ have simple normal crossings support.
Let $$L_Y:=K_Y-\lfloor \frac {B_Y}{t}\rfloor \cong \frac {1}{t}\alb _Y^*H+\{
\frac {B_Y}{t}\},$$
$$L_X:=K_X-\lfloor \frac {B_X}{t}\rfloor \cong \frac {1}{t}f^*\alb _Y^*H+\{
\frac {f^* B_Y}{t}\}.$$

\noindent {\bf Claim 1.} For $t\gg 0$, one has that $|2K_Y|=|K_Y+L_Y|
+\lfloor \frac
{B_Y}{t}\rfloor$.

\bigskip
\noindent {\bf Claim 2.} The divisor $K_{X/Y}+f^*(\lfloor B_Y/t\rfloor )
-\lfloor f^*B_Y /t \rfloor$ is effective.

\bigskip
Grant these for the time being. By Claim 2, there is an inclusion 
$$f^*(\omega _Y\ot L_Y)\lra \omega _X \ot L_X$$
and hence an inclusion
$$i: \alb _{Y,*}(\omega _Y \ot L_Y) \lra \alb _{Y,*} f_*(\omega _X \ot L_X).$$
By Theorem \ref{omegavan}, 
one has that for all $i>0$ and $P\in \Pic ^0(Y)$
$$H^i(\Alb (Y),\alb _{Y,*}(\omega _Y \ot L_Y) \ot P)=
H^i(\Alb (Y), \alb _{Y,*} f_*(\omega _X \ot L_X)\ot P)=0.$$
We also have that $$h^0(\omega _Y \ot L_Y)=P_2(Y)=P_2(X)\geq 
h^0( \omega _X \ot L_X)
\geq h^0(\omega _Y \ot L_Y)$$
and therefore for all $P\in \Pic ^0(Y)$
$$h^0(\alb _{Y,*}(\omega _Y \ot L_Y)\ot P)=\chi (\alb _{Y,*}
(\omega _Y \ot L_Y)\ot P)=$$
$$\chi (\alb _{Y,*}(\omega _Y \ot L_Y))= h^0(\alb _{Y,*}(\omega _Y \ot L_Y))=$$
$$ h^0( \alb _{Y,*}f_*(\omega _X \ot L_X))=\chi( \alb _{Y,*}
f_*(\omega _X \ot L_X))= $$
$$\chi( \alb _{Y,*}f_*(\omega _X \ot L_X)\ot P)=h^0
( \alb _{Y,*}f_*(\omega _X \ot L_X)\ot P).$$
So the inclusion $i$ induces isomorphisms in cohomology 
for all $i\geq 0$ and all
$P\in \Pic ^0(Y)$ and hence 
$$\alb _{Y,*}(\omega _Y \ot L_Y) \cong \alb _{Y,*} f_*(\omega _X \ot L_X).$$
In particular the generic degrees of $\alb _Y, \alb _Y\circ f$ coincide and
hence $f$ is birational.
We now turn to the proof of the above claims.

\bigskip \noindent{\it Proof of Claim 1.} After replacing $Y$ by an 
appropriate birational 
model, we may assume that $$|2K_Y|=|M|+F$$
where $|M|$ is free and $F+B_Y$ has simple normal crossings support.
For $l\gg 0$ let $G\in |lM|$ be general.
Consider $$B^\prime _Y:=B_Y+lF+G\in |(t+2l)K_Y-\alb _Y^*H|.$$
We may write $B_Y=\sum b_iB_i$ and $F=\sum f_iB_i$ for distinct 
reduced irreducible
divisors $B_i$ on $Y$. We have that for $l\gg 0$,
the multiplicity of $B^\prime _Y$ along $B_i$
satisfies $$\lfloor \frac {mult _{B_i}(B^\prime _Y)}
{t+2l}\rfloor =\lfloor \frac
{b_i+lf_i}{t+2l}\rfloor\leq \lfloor \frac {b_i}{2l}+\frac 
{f_i}{2}\rfloor \leq f_i
=mult _{B_i} F.$$
We now replace $t$ and $B_Y$ by $t+2l$ and $B^\prime _Y$.

\bigskip 
\noindent{\it Proof of Claim 2.} This follows from a local computation
(see also Proposition 2.8 
of \cite{Ein}). 

Let $B_Y=\sum b_iB_i$ where $B_i$ are distinct prime divisors. 
If the $b_i$ are integers, then $\lfloor (f^* B_Y)/t \rfloor =f^*
\lfloor B_Y/t \rfloor$. So we may assume that $\lfloor B_Y/t \rfloor =0$
and $B_Y/t\ne 0$.
Let $Q\in X$ be any point and $P=f(Q)\in Y$ be its image.
We may assume that $P\in B_i$ for $1\leq i\leq s$. Let $(y_1,...,y_n)$ be local
coordinates centered in $P$ such that $y_i$ is a local equation for $B_i$
for $1\leq i\leq s$. Let $E$ be any component of $f^*B_Y$ containing $Q$ and
choose local coordinates $(x_1,...,x_n)$
on $X$ centered in $Q$ such that $x_1$
is a local equation of $E$. One has that $f^*y_i=x_1^{n_i}\epsilon _i$,
with $n_i\geq 0$ and $\epsilon _i$ a regular function not vanishing on $E$.
One sees that a local equation for $K_{X/Y}$ near $Q$ is given by
the determinant of the Jacobian matrix $\left( \frac {\partial f^*y_i}
{\partial x_j} \right)$, which vanishes on $E$ to order at least
$\left( \sum n_i\right) -1$. So the coefficient of $E$ in
$K_{X/Y}-\lfloor (f^*B_Y)/t\rfloor +f^* \lfloor B_Y/t\rfloor$
is at least $$\sum n_i -\lfloor \sum n_i\frac{b_i}{t} \rfloor 
 -1\geq 0.$$
\end{proof}
We remark that the above theorem also holds if $Y$ has sufficiently large 
algebraic fundamental group (see \cite{Hac4} for a discusion of this case).  
\newpage
%%%%%%%%%%%%%%%%%%%%%%%%%%%%%%%%%%%%%%%%%%%%%%
\section{Theta divisors: Singularities and Cohomological characterization}
%%%%%%%%%%%%%%%%%%%%%%%%%%%%%%%%%%%%%%%%%%%%%%%
Let $(A,\Theta )$ be a principally polarized abelian variety
(PPAV), i.e. $\Theta$ is an ample divisor with $h^0(A,\OO (\Theta ))=1$.
The singularities of such divisors are closely related to the geometry
of the abelian variety $A$. For example, Kempf has shown that if
$A$ is the Jacobian variety of a curve, then the theta divisor has rational
singularities. More recently Koll\'ar has shown
\begin{thm} \label{K1}
{\bf (Koll\'ar)} \cite{Ko3}
Let $(A,\Theta )$ be a PPAV, then the pair 
$(A,\Theta )$ is log canonical
\footnote{In this section, we will use the formalism of multiplier ideals
and some of their applications. We refer to \cite{Ein} for an overview of
these and related concepts}. In particular if
$$\Sigma _k(\Theta )=\{ x\in A| mult _x(\Theta )\geq k \},$$
then every component of $\Sigma _k$ has codimension at least $k$. 
\end{thm}
\begin{proof}
Let $\nu :A'\lra A$ be a log resolution of the pair $(A,\Theta
)$ (i.e. a proper birational morphism such that $\nu ^{-1}(\Theta)\cup
\{ exceptional\ set\ of\ \nu \}$ is a divisor with normal crossings support).
For any fixed $1\gg \epsilon >0$ let $$\II :=\nu _*(\OO _{A'}(K_{A'/A}-
\lfloor (1-\epsilon)\nu ^* \Theta \rfloor ) \subset \OO _A$$
be the {\it multiplier ideal sheaf} of the pair $(A,\Theta)$.
Multiplier ideals are a measure of the singularities of the pair
$(A,(1-\epsilon)\Theta )$. Their definition does not depend on the choice
of the log-resolution $\nu$.
By definition, the pair $(A,\Theta )$ is log canonical if for all $1\gg
\epsilon > 0$, one has $\II ((1-\epsilon)\Theta )=\OO _A$.
We now proceed to show that $\II =\OO _A$. 
First of all, we have that
$$\nu ^* (\Theta +P) +K_{A'/A}  -
\lfloor (1-\epsilon)\nu ^* \Theta \rfloor \equiv
K_{A'}+\epsilon \nu ^* \Theta +\{(1-\epsilon) \nu ^* \Theta \}$$
and hence by Theorem \ref{omegavan}, for all $P\in \Pic ^0(A)$,
$$h^i(A, \OO _A(\Theta )\ot P \ot \II )=0\ \ \ for\ all\ i>0.$$
If $\II \ne \OO _A$, then $h^0(A, \OO _A(\Theta )\ot P \ot \II )=0$
for general $P\in \Pic ^0(A)$ (as a general translate of $\Theta$ will
not vanish on the cosupport of $\II$).
But then for all $P\in \Pic ^0(A)$,
$$h^0(A, \OO _A(\Theta )\ot P \ot \II )=
\chi (\OO _A(\Theta )\ot P \ot \II )=\chi (\OO _A(\Theta )\ot \II )=0$$
and hence (cf. Example \ref{E1}) 
$\OO _A(\Theta )\ot \II$ is trivial, 
which is impossible.
\end{proof}
The above proof, is due to Ein and Lazarsfeld. 
The same argument also shows that if $D\in |m\Theta |$ ,
then the pair $(A,(1/m)\Theta )$ is log canonical. 

One can easily construct examples of abelian varieties where the bound
on the codimension of $\Sigma _k(\Theta )$ 
is achieved by considering products of abelian varieties of smaller dimension. 
The following result of Ein and Lazarsfeld shows that this is essentially 
the only case in which such a phenomenon occurs.
\begin{thm}\label{el}
 {\bf(Ein-Lazarsfeld) \cite{EL}} If $\Theta \subset A$ is an irreducible 
theta divisor, then $\Theta$ is normal and has only rational singularities.
\end{thm}
\begin{proof} Let $\nu :A'\lra A$ be a log resolution of the pair 
$(A,\Theta )$ and $Z$ be the proper transform of $\Theta$.
Pushing forward exact sequence
$$0\lra \omega  _{A'}\lra \omega _{A'}(Z)\lra \omega _Z\lra 0$$
via $\nu :A'\lra A$, one obtains a short exact sequence\footnote{Recall that
by Theorem \ref{kollar}, $R^1\nu _* \omega _{A'}=0$.}
$$0\lra \OO _A \lra \OO _A (\Theta ) \ot \II \lra \nu_* \omega _Z\lra 0$$ 
where $\II := \nu _* (\omega _{A'/A}(Z-\nu ^* \Theta ))$ is the {\it
adjoint ideal} of $\Theta$ (cf. \cite{Ein} \S 3).
In order to show that $\Theta$ has rational singularities,
it suffices to show that $\II =\OO _A$\footnote{Since $\Theta$ is Gorenstein
and $K_{A'/A}
(Z-\nu ^* \Theta)|Z=K_Z-\nu^*K_{\Theta }$, it follows that $\II =\OO _A$
iff $\nu _* K_Z=K_{\Theta }$ i.e. iff $\Theta $ has canonical and hence
rational singularities cf. \cite{Ein}.}.
By the generic vanishing theorem Theorem \ref{GVT}, one sees that
for general $P\in \Pic ^0(A)$
$$h^i(\nu _* (\omega _Z )\ot P)=h^i(\omega _Z \ot \nu ^*P)=0\ \ for\ all\ 
i>0.$$
If $\II\ne \OO _A$, then also $h^0(\nu _*(\omega _Z) \ot P)=
h^0(\OO _A (\Theta ) \ot \II \ot P)=0$ for general $P\in \Pic ^0(A)$
and so $\chi (\nu _* \omega _Z)=0$.
This contradicts a result of Kawamata and Viehweg.  
\end{proof}
This leads to the following cohomological characterization of principal
polarizations (cf. \cite{Hac1}).
\begin{thm}\label{theta} {\bf (Hacon)} 
Let $\Theta \subset A$ be an irreducible reduced divisor of a
$g$-dimensional abelian variety,
$\nu :A'\lra A$ a log resolution of $(A,\Theta )$ and $Z$ the proper transform
of $\Theta$. Then $\Theta$ is a principal polarization
iff  $$h^i(\omega _Z)=\binom {g}{i+1}\ for \ 
0\leq i\leq g-1\ \ \ \ and$$
$$h^i(\omega _Z \ot \nu ^* P)
=0\ \forall \ \OO _A \ne P\in \Pic ^0(A),\ i>0$$
\end{thm}
\begin{proof} If $\Theta$ is a principal polarization, then
for all $P\in \Pic ^0(A)$ one has $$h^0(\OO _A(\Theta )\ot P)=1\ \ \  and
\ \ h^i(\OO _A(\Theta )\ot P)=0\ \  \forall i>0.$$
For all $P\ne \OO _A$ and $0\leq i\leq g$, one has $h^i(P)=0$.
Consider the exact sequence
$$0\lra \OO _A \lra \OO _A (\Theta )\lra \nu _*\omega _Z\lra 0.$$
Twisting it by $P\in \Pic ^0(A)$, one immediately computes
$h^i(\omega _Z \ot \nu ^* P)=h^i(\nu _*(\omega _Z )\ot P)$.

Suppose now that $h^i(\omega _Z \ot \nu ^* P)$ are as above.
By Example \ref{Ex4}, 
it suffices to show that for all $P\in \Pic ^0(A)$, one has
$h^0(\OO _A(\Theta )\ot P \ot \II )=1$ and
$h^i(\OO _A (\Theta )\ot P \ot \II )=0$ for all $i>0$
(then $\OO _A (\Theta )\ot \II$ is a line bundle and hence $\II=\OO _A$
and $h^0(\OO _A(\Theta ))=1$).
When $P\ne \OO _A$, from the above exact sequence, one sees that 
$h^i(\OO _A (\Theta )\ot P \ot \II )=h^i(\nu _*\omega _Z \ot P)$
as required. (Notice that since $h^i(\nu _*\omega _Z \ot P)=0$ for
all $i>0$, one has $$h^0(\nu _*\omega _Z \ot P)=\chi (\nu _*\omega _Z \ot P)=
\chi (\nu _*\omega _Z)=1.)$$
The homomorphism $$H^0(\Omega _A^{g-i-1})\lra 
H^0(\Omega _Z ^{g-i-1})$$ is injective.
Otherwise, there would be a morphism $\Theta \lra \bar{\Theta}=
\Theta /K$ where $K$ is a subtorous of $A$. Let $F=F_{Z/\bar{\Theta }}$,
for general $P\in \Pic ^0(A)$, one has $P|K\ne \OO _K$
and so $$h^0((\omega _Z\ot P)|F)\cong h^0(\omega _F\ot P)\cong
h^0(\omega _K \ot P)=0.$$
Therefore $h^0(Z,\omega _Z \ot P)=0$ which is impossible. 
By Hodge symmetry and Serre duality, it follows that
the homomorphisms $$H^i(\nu _* \omega _Z)\lra H^{i+1}(\OO _A)$$
is surjective. By dimensional considerations it 
is an isomorphism and hence $$h^0(\OO _A(\Theta ) \ot \II )
=h^0(\OO _A)=1\ \ \ \   and\  
\ h^i(\OO _A(\Theta ) \ot \II )=0\ for \ i>0.$$ 
\end{proof}

In general the hypothesis in the above theorem can't be weakened,
however, in \cite{HP1} it is shown that: {\em If $X$ is a smooth
projective variety of maximal Albanese dimension, with 
$P_1(X)=q(X)=\dim (X)+1$, and the Albanese variety of $X$ is simple,
then $X$ is birational to a theta divisor.}

As an application of the above result, one can also classify
surfaces with $p_g=q=3$ (cf. \cite{HP2}, \cite{Pi}). 
\newpage
%%%%%%%%%%%%%%%%%%%%%%%%%%%%%%%%%%%%%%%%%%%%%%%%%%%%
\section{Characterization of abelian varieties} 
%%%%%%%%%%%%%%%%%%%%%%%%%%%%%%%%%%%%%%%%%%%%%%%%
We begin by recalling the following result of 
Kawamata which is a generalization
to higher dimensions of Enriques' Theorem for surfaces 
mentioned in the introduction. 
\begin{thm}\label{kaw}
{\bf (Kawamata)} \cite{Ka1} Let $X$ be a smooth complex 
projective variety
with $\kappa (X)=0$, then $\alb :X\lra \Alb $ is surjective (in particular
$q(X)\leq \dim (X)$). If moreover, $q(X)=\dim (X)$, then
$X$ is birational to an abelian variety.
\end{thm}
This gives (a partial) answer to a conjecture of Ueno which in turn is
an important test case for the $C_{n,m}$ conjecture of Iitaka. On the other
hand, the invariant $\kappa (X)$ is highly nontrivial as it involves knowing
$P_m(X)$ for infinitely many values of $m$. 
Therefore it would be more satisfactory
to give an "effective" statement such as the following:
\begin{thm}\label{ko4} 
{\bf (Koll\'ar) \cite{Ko4}} Let $X$ be a smooth projective variety.
If $P_m(X)=1$ for some $m\geq 3$, then $\alb : X \lra \Alb $ is surjective.
If $P_m(X)=1$ for some $m\geq 4$ and $q(X)=\dim (X)$, then
$X$ is birational to an abelian variety.
\end{thm}
We will give an idea of the proof of the 
following result conjectured by Koll\'ar:
\begin{thm}\label{ch} {\bf (Chen-Hacon) \cite{CH1}, \cite{CH2}}
Let $X$ be a smooth projective variety.
\begin{enumerate}
\item If $P_m(X)=1$ for some $m\geq 2$, then $\alb : 
X \lra \Alb $ is surjective.
\item If $P_m(X)=1$ for some $m\geq 2$ and $q(X)=\dim (X)$,
then $X$ is birational to an abelian variety.
\end{enumerate}
\end{thm}
\begin{proof} 
If $\kappa (X)=0$, both results follow from Kawamata's Theorem above.
We will therefore assume that $\kappa (X)>0$ and derive a contradiction
(i.e. that $P_2(X)\geq 2$).

(1) Let $f:X\lra V$ be a birational model of the Iitaka fibration of 
$X$\footnote{ By definition, one has that $\kappa (X_v)=0$, $f$ is surjective,
the general fiber is irreducible and $\dim (V)=\kappa (X)$. 
For all sufficiently big and divisible integers $m$, 
the morphism $f$ is birational
to $\varphi _m$.}.
Let $X_v$ be a generic geometric fiber of $f$, then $\kappa (X_v)=0$
and by Kawamata's Theorem, one sees that $\alb (X_v)$ is an 
abelian subvariety of $\Alb$. Since abelian subvarieties of 
a fixed abelian variety
$A$ are rigid, it follows that there is an abelian subvariety $K\subset \Alb$
such that for general $v$, the images $\alb (X_v )$ are translates of $K$.
Let $S=A/K$. We will denote by $Z,W$ the images of $X,V$ in $A,S$.
After replacing $X,V$ by appropriate birational models,
we may assume that there is a commutative diagram:
$$
\begin{CD}
X @>>> Z @>{\subset}>>A \\
@V{f}VV  & &     @V{\hat{f}}VV \\
V @>>> W @>{\subset}>> S
\end{CD}
$$
Let $\pi :X \lra S$ be the induced morphism.
Fix $H$ an ample divisor on $S$. For $r\gg 0$ the linear series
$|rK_X-\pi ^* H|$ will be non empty. After replacing $X$
by an appropriate birational model, we may assume that
$$|rK_X-\pi ^* H|=F+|M|$$
with $|M|$ base point free and $F$ with normal crossings support.
Fix $B$ a general member in $|rK_X-\pi ^*H|$ and let
$$L:=\omega _X \ot \OO _X\left( -\lfloor \frac {B}{r}\rfloor \right)
\equiv \frac {\pi ^*H}{r}-\{ \frac {B}{r} \} .$$
(One may think of $L$ as being the positive part of $\omega _X$
and $ \lfloor \frac {B}{r}\rfloor$ as being some negligible divisor.
For example if $K_X$ is ample, take $L=\omega _X$ and 
$ \lfloor \frac {B}{r}\rfloor =0$.)
As in the proof of Theorem \ref{birational}, 
one may arrange that $$|2K_X|=|K_X +L|+\lfloor
B/r\rfloor .$$ Since $ {H}/{r}$ is ample and $\{ B/r\}$ is KLT
(it has simple normal crossings support and $\lfloor \{ B/r\}
\rfloor =0$), by Theorem \ref{omegavan}
one has that $$h^i(W, \pi _* (\omega _X \ot L)\ot P)=0\ \ \  for\  all\ i>0
\ \ \ and\ all\ P\in \Pic ^0(S).$$
Let $$\FF:=\pi _* (\omega _X \ot L).$$
One has
$$h^0(W,\FF )=h^0( W, \pi _* (\omega _X \ot L))=h^0(X, \omega _X \ot L)=
h^0(X,\omega _X ^{\ot 2})$$
It follows that 
$$h^0(W, \FF \ot P)=\chi (W, \FF \ot P)=\chi (W, \FF )=h^0(W,\FF )=1.$$
By Example \ref{Ex4}, any sheaf $\FF$ as 
above must be supported on a translate
of an abelian subvariety of $S$. By Theorem \ref{kollar}, 
the sheaf $\FF$ is a torsion free
sheaf on $W$. However, $W$ generates $S$ and hence $W=S$
and so $Z=\Alb$ i.e. $\alb :X\lra \Alb$ is surjective.

(2) Here we will only consider the case $P_m(X)=1$ for some $m\geq 4$.
By the proof of
part (1), one sees that $$|K_X+L+\eta|>0\ \ \  for\ all\  \eta \in 
\Pic ^0(S).$$
From the morphism $$|K_X+L+\eta|\times |K_X+L-\eta|\lra
|2K_X+2L|$$
one sees that
$$P_4(X)\geq h^0(\omega _X ^{\ot 2}\ot L^{\ot 2})\geq \dim (S)+1$$
and so $\kappa (V)=\dim (V)\leq \dim (S)=0$ and the result follows 
from Kawamata's 
theorem mentioned above.
\end{proof}


We now give an idea of the proof in the case $m=2$.
Assume that $\kappa (X)=\dim (X)$. By a result of Ein and Lazarsfeld 
(see \cite{CH1}),
one has that if $\kappa (X)>0$, then
$$\dim V^0(X,\omega _X)>0.$$
So fix $T$ a positive dimensional component of $V^0(X,\omega _X)$.
By Simpson's result (cf. Theorem \ref{GVT}), we may write
$$T=P_0+T_0$$ where $P_0$ is a torsion element
of $\Pic ^0(X)$ and $T_0$ is a sub-abelian variety of $\Pic ^0(X)$.
Consider now for all $Q\in T_0$ the map
$$|K_X+P_0+Q|\times |K_X +P_0-Q|\lra |2K_X+2P_0|.$$
Since for all $Q\in T_0$, the linear series $|K_X+P_0+Q|,|K_X +P_0-Q|$
are non-empty and any effective divisor can be decomposed in to the 
sum of two effective divisors in finitely many ways, one has that 
$$h^0(X,\omega _X ^{\ot 2}\ot P_0^{\ot 2})\geq 2.$$
Proceeding as in the proof of Theorem \ref{birational},
one sees that $$h^0(\omega _X ^{\ot 2}\ot P_0^{\ot 2})= P_2(X).$$
This is the required contradiction.

The case when $0<\kappa (X) <\dim (X)$ is based on a similar idea,
but it is unluckily much more technical.
We refer the reader to \cite{CH1}
for this part of the proof.



\begin{thebibliography}{Ka12}
\bibitem[CH1]{CH1}
{\sc Chen J.\,A.} and {\sc Hacon C\,D.}:
Characterization of abelian varieties,
Invent. Math. {\bf 143}, 435-447 (2001)

\bibitem[CH2]{CH2}
{\sc Chen, J.\,A.} and {\sc Hacon, C.\,D.}:
On algebraic fiber spaces over varieties of maximal albanese dimension,
Duke Math. J., Vol. 111, No. 1, 2002

\bibitem[CH3]{CH3}
{\sc Chen, J.\,A.} and {\sc Hacon, C.\,D.}:
{Linear series of irregular varieties},
Proceedings  of the symposium
on Algebraic Geometry in East Asia. World Scientific (2002)

\bibitem[CH4]{CH4}
{\sc Chen, J.\,A.} and {\sc Hacon, C.\,D.}:
On the irregularity of the image of the Iitaka fibration.
To appear in Comm. in Algebra.

\bibitem[Ein]{Ein}
{\sc Ein, L.}:
Multiplier ideals, vanishing theorems and applications,
Algebraic Geometry (Santa Cruz 1995), Proc. Sympos. 
Pure Math. {\bf 62} (1), Amer. Math. Soc.,
Providence, RI, 1997, pp. 203-219

\bibitem[EL]{EL}
{\sc Ein, L.} and {\sc Lazarsfeld, R.}:
Singularities of theta divisors, and the birational geometry of irregular
varieties, J. Amer. Math. Soc {\bf 10} (1997), 243-258

%\bibitem[EV]{EV}
%{\sc Esnault, H.} and {\sc Viehweg W.}:
%Lectures on Vanishing Theorems, DMV Seminar, {\bf 20}. 
%Birkh\"{a}user Verlag, 1992

%\bibitem[Fu]{Fu}
%{\sc Fujino, O.}:
%Algebraic fiber spaces whose general  fibers
%are of maximal albanese dimension. Preprint math.AG/0204262


\bibitem[GL]{GL}
{\sc Green, M.} and {\sc Lazarsfeld, R.}:
Higher obstructions to deforming cohomology groups of line bundles,
Jour. Amer. Math. Soc.
vol. 4, {\bf 1}, 1991,  87-103

\bibitem[Hac1]{Hac1}
{\sc Hacon, C.\,D.}: Fourier transforms, generic vanishing theorems
and polarizations of abelian varieties, Math. Zeitschrift {\bf 235} (2000)
717-726

\bibitem[Hac2]{Hac2}
{\sc Hacon, C.\,D.}:
Varieties with $P_3=3$ and $q=\dim (X)$, to appear in Math. Nach.

\bibitem[Hac3]{Hac3}
{\sc Hacon, C.\,D.}:
Effective criteria for birational morphisms.
Jour. London Math. Soc. Vol. {\bf 67} part 2 April 2003, 337-348

\bibitem[Hac4]{Hac4}
{\sc Hacon, C.\,D.}:
A derived category approach to generic vanishing theorems.
Preprint

\bibitem[HP1]{HP1}
{\sc Hacon, C.\,D.} and {\sc Pardini, R.}:
On the birational geometry of varieties of
maximal Albanese dimension, Jour. Reine Angew. Math. {\bf 546} (2002)

\bibitem [HP2]{HP2} C. D. Hacon and R. Pardini,
{\em Surfaces with $p_g=q=3$},
Trans. of the Amer. Math. Soc. {\bf 354}, No. 7, 2631-2638

\bibitem[Ha1]{Ha1}
{\sc Hartshorne, R.}:
Algebraic Geometry, Graduate texts in Math., 1977, {\bf 52} Springer Verlag

%\bibitem[Ha2]{Ha2}
%{\sc Hartshorne, R.}:
%Residues and Duality, L.N.M. {\bf 20}, 1996, Springer Verlag

\bibitem[Ka1]{Ka1}
{\sc Kawamata, Y.}:
Characterization of abelian varieties,
Compositio Math. {\bf 43}, 1981, 253-276


\bibitem[Ko1]{Ko1}
{\sc Koll\'ar, J.}:
Higher direct images of dualizing sheaves I,
 Ann. Math. {\bf 123},
1986, 11--42

\bibitem[Ko2]{Ko2}
{\sc Koll\'ar, J.}:
Higher direct images of dualizing sheaves II,
 Ann. Math. {\bf 124},
1986, 171--202

\bibitem[Ko3]{Ko3}
{\sc Koll\'ar, J.}:
Shafarevich maps and automorphic forms,
M. B. Porter Lectures, Princeton Univ. Press, Princeton, 1995

\bibitem[Ko4]{Ko4}
{\sc Koll\'ar, J.}:
Shafarevic maps and plurigenera of algebraic varieties, Invent. Math {\bf 113},
1993, 177-215

%\bibitem[La]{La}
%{\sc Lazarsfeld, R.}:
%Positivity in Algebraic Geometry,
%Preprint

%\bibitem[Mo]{Mo}
%{\sc Mori, S.}:
%Classification of higher dimensional varieties, Algebraic Geometry 
%(Brunswick, Maine, 1985), Proc. Sympos. Pure Math. {\bf 46},
%Amer. Math. Soc., Providence, 1987, 269-331

\bibitem[Mu]{Mu}
{\sc Mukai, S.}:
Duality between $D(X)$ and $D(\hat {X})$
with its application to Picard sheaves,
Nagoya math. J. {\bf 81},
1981, 153-175

\bibitem[Pi]{Pi}
G. P. Pirola,
{\em Surfaces with $p_g=q=3$}, Manuscripta Mathematica (to appear).

%\bibitem[V1]{V1}
%{\sc Viehweg, E.}:
%Weak positivity and the additivity of the Kodaira dimensionof certain fiber spaces.
%Adv. Studies in Pure Math 1, 1983, 329-353

%\bibitem[V2]{V2}
%{\sc Viehweg, E.}:
%Quasi-projective moduli for polarized manifolds,
%Ergenbisse der Mathematik 30 (Springer, 1995)


\end{thebibliography}

\end{document}
\end

